\documentclass[11pt]{article}
\usepackage[margin=1in]{geometry}
\usepackage{amsmath}
\usepackage{amssymb}
\usepackage{tcolorbox}

\title{Problem 2c: Invalid Solution \\ (Why the Wrong Guess Fails)}
\author{ECO 310 - Learning from Mistakes}
\date{February 23, 2026}

\begin{document}

\maketitle

\section{What Problem 2c Asks}

\textbf{Question:} Show that you get an invalid solution to (b) if $p_1 < p_2$. And explain why intuitively, the guess was wrong about which nonnegativity constraint binds.

\textbf{Translation:} Use the SAME guess from part (b), but now assume $p_1 < p_2$ (good 1 is cheaper). Show that this leads to a contradiction.

\section{Recall: The Guess from Part 2b}

The guess was:
\begin{itemize}
\item Utility constraint is binding: $x_1 + x_2 = u^*$
\item Nonnegativity on good 1 is binding: $x_1 = 0$
\item Nonnegativity on good 2 is non-binding: $x_2 > 0$
\end{itemize}

This means: buy only good 2, so $x_1 = 0$ and $x_2 = u^*$.

\section{Apply the Same Logic with $p_1 < p_2$}

Let's follow the exact same steps as part (b), but now assume $p_1 < p_2$ (good 1 is CHEAPER).

\subsection{Step 1: Use Complementary Slackness}

Since $x_2 > 0$:
\begin{equation}
\lambda_2 = 0
\end{equation}

\subsection{Step 2: Solve for $\lambda_3$}

From the FOC for $x_2$ (which must equal zero since $x_2 > 0$):
\begin{align}
-p_2 + \lambda_2 + \lambda_3 &= 0 \\
-p_2 + 0 + \lambda_3 &= 0 \\
\lambda_3 &= p_2
\end{align}

\subsection{Step 3: Solve for $\lambda_1$}

From the FOC for $x_1$:
\begin{equation}
-p_1 + \lambda_1 + \lambda_3 \leq 0
\end{equation}

Substitute $\lambda_3 = p_2$:
\begin{align}
-p_1 + \lambda_1 + p_2 &\leq 0 \\
\lambda_1 &\leq p_1 - p_2
\end{align}

At the optimum, assume this binds:
\begin{equation}
\lambda_1 = p_1 - p_2
\end{equation}

\subsection{Step 4: Check Validity}

Now here's the problem! We assumed $p_1 < p_2$, so:
\begin{align}
\lambda_1 &= p_1 - p_2 \\
&< 0 \quad \text{(since $p_1 < p_2$)}
\end{align}

\begin{tcolorbox}[colback=red!10!white,colframe=red!75!black,title=INVALID!]
\textbf{Problem:} We got $\lambda_1 < 0$, which violates the KKT condition that all Lagrange multipliers must be non-negative!

$$\lambda_1 \geq 0 \text{ is required, but we got } \lambda_1 = p_1 - p_2 < 0$$

This means the guess $x_1 = 0$, $x_2 = u^*$ is NOT optimal when $p_1 < p_2$.
\end{tcolorbox}

\section{Why Is This Invalid?}

\begin{tcolorbox}[colback=yellow!10!white,colframe=orange!75!black,title=The Rule]
\textbf{KKT Condition:} All Lagrange multipliers must satisfy $\lambda_i \geq 0$.

If any $\lambda_i < 0$, the solution is not optimal.
\end{tcolorbox}

In our case:
\begin{itemize}
\item We found $\lambda_1 = p_1 - p_2$
\item When $p_1 < p_2$, this gives $\lambda_1 < 0$
\item This violates the non-negativity requirement
\item Therefore, the guess is invalid
\end{itemize}

\section{Intuitive Explanation}

\subsection{What Does $\lambda_1 < 0$ Mean?}

Remember: $\lambda_1$ is the shadow price of the constraint $x_1 \geq 0$.

\begin{itemize}
\item $\lambda_1 > 0$ means: "I'd benefit from relaxing $x_1 \geq 0$ (buying negative $x_1$)"
\item $\lambda_1 = 0$ means: "The constraint $x_1 \geq 0$ doesn't limit me"
\item $\lambda_1 < 0$ means: "Relaxing $x_1 \geq 0$ would HURT me" - but this is nonsense!
\end{itemize}

\textbf{The contradiction:} If $\lambda_1 < 0$, it means you'd be WORSE OFF if you could buy $x_1 < 0$ (negative good 1). But you're already at $x_1 = 0$, so this constraint isn't actually binding you in a meaningful way.

\subsection{The Economic Intuition}

\begin{tcolorbox}[colback=green!5!white,colframe=green!75!black,title=Why the Guess Was Wrong]
\textbf{The guess:} Buy only good 2 ($x_1 = 0$, $x_2 = u^*$)

\textbf{The reality:} When $p_1 < p_2$ (good 1 is CHEAPER), you should buy only good 1, not good 2!

\textbf{What went wrong:} We guessed that the consumer buys the MORE EXPENSIVE good, which is obviously not optimal for cost minimization.
\end{tcolorbox}

\subsection{Concrete Example}

Suppose:
\begin{itemize}
\item Coffee costs \$2 per cup ($p_1 = 2$)
\item Tea costs \$5 per cup ($p_2 = 5$)
\item You need utility 10 ($u^* = 10$)
\end{itemize}

\textbf{The wrong guess:} Buy 0 coffee, 10 teas
\begin{itemize}
\item Cost: $2(0) + 5(10) = \$50$
\end{itemize}

\textbf{The right choice:} Buy 10 coffees, 0 teas
\begin{itemize}
\item Cost: $2(10) + 5(0) = \$20$
\end{itemize}

Obviously, buying only the MORE EXPENSIVE good (\$50) is worse than buying only the CHEAPER good (\$20)!

When we apply the math:
\begin{align}
\lambda_1 &= p_1 - p_2 \\
&= 2 - 5 \\
&= -3 < 0 \quad \text{INVALID!}
\end{align}

\section{What Should the Correct Guess Be?}

When $p_1 < p_2$ (good 1 is cheaper), the correct guess should be:
\begin{itemize}
\item $x_1 = u^*$ (buy only good 1)
\item $x_2 = 0$ (buy no good 2)
\item Nonnegativity on good 1 is NON-binding ($x_1 > 0$ so $\lambda_1 = 0$)
\item Nonnegativity on good 2 IS binding ($x_2 = 0$ so $\lambda_2 \geq 0$)
\end{itemize}

This is the OPPOSITE of what we guessed in part (b)!

\section{Complete Answer to Problem 2c}

\begin{tcolorbox}[colback=blue!5!white,colframe=blue!75!black,title=Problem 2c Answer]

\textbf{Showing the solution is invalid:}

Using the guess from part (b) with $p_1 < p_2$, we get:
\begin{align}
x_1 &= 0 \\
x_2 &= u^* \\
\lambda_1 &= p_1 - p_2 < 0 \quad \text{(INVALID - violates $\lambda_1 \geq 0$)} \\
\lambda_2 &= 0 \\
\lambda_3 &= p_2
\end{align}

Since $\lambda_1 < 0$, this violates the KKT conditions. The solution is invalid.

\textbf{Intuitive explanation:}

The guess assumed we buy only good 2 ($x_1 = 0$, $x_2 = u^*$). But when $p_1 < p_2$, good 1 is CHEAPER, so buying only the more expensive good 2 is clearly not cost-minimizing.

The negative $\lambda_1$ tells us mathematically what we know intuitively: we should be buying MORE of good 1 (the cheaper good), not zero.

The correct guess when $p_1 < p_2$ should have nonnegativity on good 2 binding (buy only good 1), not nonnegativity on good 1 binding.

\end{tcolorbox}

\section{General Pattern for Perfect Substitutes}

\begin{tcolorbox}[colback=yellow!10!white,colframe=orange!75!black,title=The Rule]
For perfect substitutes:
\begin{itemize}
\item If $p_1 < p_2$: Buy only good 1 ($x_1 = u^*$, $x_2 = 0$)
\item If $p_1 > p_2$: Buy only good 2 ($x_1 = 0$, $x_2 = u^*$)
\item If $p_1 = p_2$: Any combination with $x_1 + x_2 = u^*$ works
\end{itemize}

The KKT conditions will only be satisfied (all $\lambda_i \geq 0$) when you guess the RIGHT corner solution.
\end{tcolorbox}

\section{Summary}

\begin{tcolorbox}[colback=blue!5!white,colframe=blue!75!black,title=Key Takeaway]
\textbf{The point of Problem 2c:} Show that the Lagrangian method catches our mistakes!

If you guess the wrong corner solution:
\begin{enumerate}
\item The math will produce a negative Lagrange multiplier
\item This violates the KKT conditions
\item This tells you your guess was wrong
\item You need to try a different guess
\end{enumerate}

This is how you systematically check all possible corner solutions until you find the right one.
\end{tcolorbox}

\end{document}
